\documentclass[10pt,a4paper]{article}
\usepackage[T1]{fontenc}
\usepackage[utf8]{inputenc}
\usepackage{tgtermes,tgheros,microtype}
\usepackage[a4paper,margin=18mm,headheight=19pt]{geometry}
\usepackage{xcolor,graphicx,booktabs,tabularx,array,enumitem}
\usepackage[most]{tcolorbox}
\usepackage{fancyhdr,lastpage,titlesec}
\usepackage[hidelinks]{hyperref}
\usepackage{xurl,caption,placeins,textcomp,amsmath}

\definecolor{Navy}{HTML}{173753}
\definecolor{Teal}{HTML}{008F8C}
\definecolor{PaleTeal}{HTML}{E5F4F3}
\definecolor{Orange}{HTML}{F08A24}
\definecolor{PaleOrange}{HTML}{FFF1E2}
\definecolor{Red}{HTML}{BB3E35}
\definecolor{PaleRed}{HTML}{F9E9E7}
\definecolor{Slate}{HTML}{5B7083}
\definecolor{Light}{HTML}{F3F6F8}
\definecolor{Line}{HTML}{D7E0E7}
\definecolor{Green}{HTML}{2E8B57}
\definecolor{PaleGreen}{HTML}{EAF5EE}

\newcommand{\sans}{\fontfamily{qhv}\selectfont}
\newcommand{\reportdate}{1 September 2026}
\newcommand{\reportversion}{Version 1}
\newcolumntype{Y}{>{\raggedright\arraybackslash}X}
\newcolumntype{R}{>{\raggedleft\arraybackslash}X}
\newcolumntype{C}{>{\centering\arraybackslash}X}

\hypersetup{
  pdftitle={HeiQ XMembrane EO Vapour Experiment Design -- Moisture Target and Cell Scaling},
  pdfauthor={HeiQ XMembrane project team},
  pdfsubject={Coupon-cell sizing, moisture transport target, sensitivity and reversible cycling},
  pdfkeywords={XMembrane, electro-osmosis, ERV, vapour cell, moisture flux, transport number}
}

\pagestyle{fancy}
\fancyhf{}
\fancyhead[L]{\sans\footnotesize\color{Slate}HEIQ MATERIALS AG · XMEMBRANE}
\fancyhead[R]{\sans\footnotesize\color{Slate}EO VAPOUR EXPERIMENT · V1}
\fancyfoot[L]{\sans\scriptsize\color{Slate}Experiment-design basis · conditional model evidence, not material validation}
\fancyfoot[R]{\sans\scriptsize\color{Slate}Page \thepage\ of \pageref*{LastPage}}
\renewcommand{\headrulewidth}{0.45pt}
\renewcommand{\headrule}{\hbox to\headwidth{\color{Navy}\leaders\hrule height \headrulewidth\hfill}}
\renewcommand{\footrulewidth}{0pt}

\titleformat{\section}{\sans\bfseries\color{Navy}\fontsize{19}{22}\selectfont}{\thesection}{0.65em}{}
\titleformat{\subsection}{\sans\bfseries\color{Navy}\fontsize{13}{16}\selectfont}{\thesubsection}{0.55em}{}
\titlespacing*{\section}{0pt}{7mm}{3mm}
\titlespacing*{\subsection}{0pt}{5mm}{2mm}
\setlength{\parindent}{0pt}
\setlength{\parskip}{1.6mm}
\setlist[itemize,enumerate]{leftmargin=5mm,itemsep=1mm,topsep=1mm}
\captionsetup{font=small,labelfont={bf,color=Navy}}
\renewcommand{\arraystretch}{1.18}

\newtcolorbox{mainbox}{enhanced,breakable,colback=PaleTeal,colframe=Teal,boxrule=0.7pt,arc=2mm,left=4mm,right=4mm,top=3mm,bottom=3mm}
\newtcolorbox{warningbox}{enhanced,breakable,colback=PaleOrange,colframe=Orange,boxrule=0.7pt,arc=2mm,left=4mm,right=4mm,top=3mm,bottom=3mm}
\newtcolorbox{stopbox}{enhanced,breakable,colback=PaleRed,colframe=Red,boxrule=0.7pt,arc=2mm,left=4mm,right=4mm,top=3mm,bottom=3mm}
\newtcolorbox{gobox}{enhanced,breakable,colback=PaleGreen,colframe=Green,boxrule=0.7pt,arc=2mm,left=4mm,right=4mm,top=3mm,bottom=3mm}
\newtcolorbox{plainbox}{enhanced,breakable,colback=Light,colframe=Line,boxrule=0.5pt,arc=2mm,left=4mm,right=4mm,top=3mm,bottom=3mm}
\newcommand{\eyebrow}[1]{{\sans\bfseries\color{Teal}\MakeUppercase{#1}}\par\vspace{1.5mm}}
\newcommand{\metric}[3]{\begin{minipage}[t]{0.30\textwidth}{\sans\bfseries\fontsize{18}{20}\selectfont\color{Navy}#1}\par{\sans\bfseries\small #2}\par{\scriptsize\color{Slate}#3}\end{minipage}}

\begin{document}
\thispagestyle{empty}
\vspace*{7mm}
\begin{tcolorbox}[enhanced,colback=Navy,colframe=Navy,boxrule=0pt,arc=0mm,left=12mm,right=12mm,top=12mm,bottom=12mm]
{\sans\bfseries\large\color{white}HEIQ XMEMBRANE PROJECT TEAM}\par
\vspace{8mm}
{\sans\bfseries\fontsize{27}{31}\selectfont\color{white}EO vapour experiment design}\par
\vspace{2mm}
{\sans\fontsize{17}{20}\selectfont\color{white}Moisture target, coupon-cell scaling and reversible cycling}\par
\vspace{7mm}
{\sans\large\color{white}\reportdate\quad · \quad\reportversion}
\end{tcolorbox}

\vspace{10mm}
{\sans\bfseries\Large\color{Navy}Experiment context}\par
\vspace{4mm}
This report translates the selected O30 building-scale ERV+EO result into a
measurable coupon experiment. EO is located inside the ERV membrane/core. The
report explains the moisture-rate target, active-area selection, channel and
fixture dimensions, detectability basis, transient closed-cell response and
the limitations of symmetric water shuttling without through-flow.

\begin{mainbox}
\textbf{Primary decision.} Build one modular fixture around a 50~$\times$~50~mm
maximum active aperture (25~cm\textsuperscript{2}). Use controlled once-through
gas to determine sustained polarity-dependent EO transport. Use a sealed or
recirculating configuration only as a separate transient experiment: symmetric
polarity cycling moves water in both directions but gives zero net EO transport
over a complete cycle unless a rectifying or water-removal mechanism is added.
\end{mainbox}

\vfill
\begin{tabularx}{\textwidth}{@{}lY@{}}
\textbf{System-derived net EO target} & 200--250~g/m\textsuperscript{2}/h below 1~V \\
\textbf{Maximum active aperture} & 25~cm\textsuperscript{2}; physical coupon approximately 80~$\times$~80~mm \\
\textbf{Nominal gas state} & 32~$^\circ$C, 101,325~Pa, water activity 0.50 \\
\textbf{Nominal channel} & 2~mm gap per side; 1/2/4-mm validation sweep \\
\textbf{Nominal gas flow} & 6~L/min per side for the declared detectability screen \\
\textbf{Claim status} & Experiment-design target; no measured membrane or commercial-system validation \\
\end{tabularx}

\clearpage
\section*{Executive summary}
\addcontentsline{toc}{section}{Executive summary}

The O30 system coordinate requires 4,365.9~kg/h total water transfer through
the active ERV membrane. The matched passive ERV provides approximately
3,137.6~kg/h, leaving 1,228.3~kg/h as the additional duty associated with EO.
Dividing that incremental duty by the selected 5,899~m\textsuperscript{2}
active area gives 208~g/m\textsuperscript{2}/h. This is a system-equivalent
average, not an intrinsic membrane constant or a uniform local flux.

For the 25-cm\textsuperscript{2} coupon, 208~g/m\textsuperscript{2}/h equals
0.520~g/h or 0.144~mg/s. The ideal gross mapping at transport number
$t_w=50$ and current density $j=10$~A/m\textsuperscript{2} gives
336~g/m\textsuperscript{2}/h, corresponding to 0.840~g/h or 0.233~mg/s on
the same coupon. Gross flux is not net flux: passive back-diffusion, gas films,
storage, leakage and electrode chemistry must be resolved experimentally.

The 25-cm\textsuperscript{2} aperture is retained because it gives the largest
absolute signal and lowest perimeter-to-area ratio in the existing fixture.
At 6~L/min per side, the ideal inverse-design anchor needs 13.09~cm\textsuperscript{2}
at 0.5~A/m\textsuperscript{2} and 6.55~cm\textsuperscript{2} at
1~A/m\textsuperscript{2} to exceed the declared five-to-one humidity-signal
criterion. Order-unity transport numbers require lower flow or an integrating
water-rate measurement rather than a larger permanent fixture.

\begin{center}
\metric{208 g/m\textsuperscript{2}/h}{System-equivalent EO increment}{Use 200--250 as the coupon target}\hfill
\metric{0.520 g/h}{25-cm\textsuperscript{2} target rate}{0.144 mg/s at 208 g/m\textsuperscript{2}/h}\hfill
\metric{0 net/cycle}{Symmetric sealed cycling}{Water shuttles, then returns}
\end{center}

\begin{warningbox}
\textbf{Interpretation boundary.} A closed-cell oscillation can demonstrate
directionality, reversibility and finite water displacement. It cannot by
itself demonstrate sustained ERV dehumidification. The internal gas inventory
is small, the humidity gradient develops rapidly, blocking electrodes charge,
and a symmetric full cycle returns the EO-displaced water.
\end{warningbox}

\clearpage
\tableofcontents
\clearpage

\section{From building moisture duty to a coupon target}

For dry-air mass flow $\dot m_{da}$ and humidity ratios $w$, the required water
rate is
\[
\dot m_w=\dot m_{da}(w_{\mathrm{in}}-w_{\mathrm{target}})3600.
\]

\begin{table}[htbp]
\centering
\caption{O30 moisture balance used to define the experimental target.}
\begin{tabularx}{\textwidth}{@{}Y R Y@{}}
\toprule
Quantity & Value & Meaning \\
\midrule
Total active-ERV membrane transport & 4,365.9 kg/h & Passive plus electrically enhanced transfer \\
Matched passive-ERV transport & $\approx$3,137.6 kg/h & Same ERV with EO switched off \\
Additional duty associated with EO & 1,228.3 kg/h & Difference between total and passive cases \\
Selected active membrane area & 5,899 m\textsuperscript{2} & Coupled system-sweep result \\
Average total membrane flux & 740 g/m\textsuperscript{2}/h & Total transport divided by area \\
Area-normalised EO increment & 208 g/m\textsuperscript{2}/h & Additional EO duty divided by area \\
\bottomrule
\end{tabularx}
\end{table}

\[
J_{\mathrm{total}}=\frac{4{,}365.9}{5{,}899}
=0.740~\mathrm{kg\,m^{-2}\,h^{-1}},
\qquad
J_{\mathrm{EO,eq}}=\frac{1{,}228.3}{5{,}899}
=0.208~\mathrm{kg\,m^{-2}\,h^{-1}}.
\]

The 5,899-m\textsuperscript{2} area is generated by a coupled ERV model with
axially varying activity, temperature, current, conductivity, gas films and
back-diffusion. It is not obtained by dividing the water duty by a single
measured membrane flux. Consequently, the experimental requirement is stated
as a band: demonstrate 200--250~g/m\textsuperscript{2}/h of repeatable net
polarity-odd transport below 1~V under representative gas--gas conditions.

\section{Transport number, current and sensitivity}

The ideal gross EO relationship is
\[
J_{\mathrm{EO,gross}}=t_w\frac{j}{F}M_w
\approx0.672\,t_wj\quad\mathrm{g\,m^{-2}\,h^{-1}},
\]
where $t_w$ is mol water per mol elementary charge and $j$ is in
A/m\textsuperscript{2}. This is a gross Faradaic mapping; it does not subtract
passive diffusion, films, leakage, water storage or reaction water.

\begin{table}[htbp]
\centering
\caption{Gross EO sensitivity at 10 A/m\textsuperscript{2}.}
\begin{tabular}{@{}rr@{}}
\toprule
$t_w$ & Gross EO flux [g/m\textsuperscript{2}/h] \\
\midrule
10 & 67 \\
20 & 134 \\
30 & 202 \\
50 & 336 \\
75 & 504 \\
100 & 672 \\
\bottomrule
\end{tabular}
\end{table}

The ideal 208-g/m\textsuperscript{2}/h target corresponds to $t_wj\approx310$.
Examples are $t_w=50$ at 6.2~A/m\textsuperscript{2}, $t_w=30$ at
10.3~A/m\textsuperscript{2}, and $t_w=20$ at 15.5~A/m\textsuperscript{2}.
A real experiment can require a higher product because the observable net flux
is the gross EO contribution minus opposing transport.

\begin{plainbox}
\textbf{First-order sensitivities.} Gross flux scales linearly with $t_w$ and
$j$. Total rate scales linearly with active area at unchanged local conditions.
Voltage, heating, concentration polarisation, conductivity and back-diffusion
need not scale linearly. A high liquid-state $t_w$ is therefore a screening
lead, not a validated vapour-state coefficient.
\end{plainbox}

\section{Cell area and measurable water rate}

For uniform flux $J$ and active area $A$, $\dot m=JA$. Table~\ref{tab:rates}
compares the system-equivalent target with the ideal $t_w=50$, $j=10$ gross
reference.

\begin{table}[htbp]
\centering
\caption{Coupon rate as a function of active aperture.\label{tab:rates}}
\begin{tabularx}{\textwidth}{@{}Y R R R@{}}
\toprule
Aperture & Area & Rate at 208 g/m\textsuperscript{2}/h & Rate at 336 g/m\textsuperscript{2}/h \\
\midrule
20 $\times$ 20 mm & 4 cm\textsuperscript{2} & 0.083 g/h; 0.023 mg/s & 0.134 g/h; 0.037 mg/s \\
30 $\times$ 30 mm & 9 cm\textsuperscript{2} & 0.187 g/h; 0.052 mg/s & 0.302 g/h; 0.084 mg/s \\
40 $\times$ 40 mm & 16 cm\textsuperscript{2} & 0.333 g/h; 0.092 mg/s & 0.538 g/h; 0.149 mg/s \\
50 $\times$ 50 mm & 25 cm\textsuperscript{2} & 0.520 g/h; 0.144 mg/s & 0.840 g/h; 0.233 mg/s \\
\bottomrule
\end{tabularx}
\end{table}

The aperture sweep should be fitted to
\[
\dot m_{\mathrm{measured}}=J_AA+k_PP+\dot m_{\mathrm{blank}},
\]
where $P$ is perimeter. The 4-cm\textsuperscript{2} square and 8~$\times$~50-mm
rectangle have equal area but different perimeter and flow orientation. They
help distinguish membrane-area transport from seal, edge and distribution
artefacts.

\section{Detectability basis for the 25-cm\textsuperscript{2} fixture}

The declared sizing state is 32~$^\circ$C, 101,325~Pa, inlet water activity
0.50 and 6~L/min per side. Each humidity channel has 0.1 RH percentage point
expanded uncertainty. For two independent channels,
\[
u_{\Delta RH}=\sqrt{0.1^2+0.1^2}=0.141,
\qquad
\Delta RH_{\mathrm{required}}=5u_{\Delta RH}=0.707
\text{ percentage point}.
\]

At fixed flow, the ideal outlet signal scales approximately as
$t_wjA/\dot V$. The sizing calculation converts gross water rate to humidity
ratio and RH at the declared state, then solves for the area that reaches the
0.707-point threshold.

\begin{table}[htbp]
\centering
\caption{Minimum ideal active area at 6 L/min per side.}
\begin{tabularx}{\textwidth}{@{}Y R Y@{}}
\toprule
Gross-signal case & Minimum area & Smallest insert \\
\midrule
Inverse-design anchor, 0.5 A/m\textsuperscript{2} & 13.09 cm\textsuperscript{2} & 16 cm\textsuperscript{2} \\
Inverse-design anchor, 1.0 A/m\textsuperscript{2} & 6.55 cm\textsuperscript{2} & 9 cm\textsuperscript{2} \\
Selected ERV profile minimum, 1.242 A/m\textsuperscript{2} & 5.27 cm\textsuperscript{2} & 9 cm\textsuperscript{2} \\
Selected ERV profile maximum, 3.366 A/m\textsuperscript{2} & 1.94 cm\textsuperscript{2} & 4 cm\textsuperscript{2} \\
Order-unity $t_w=1$, 10 A/m\textsuperscript{2} & 65.5 cm\textsuperscript{2} & Does not fit \\
Low-drag $t_w=0.3$, 10 A/m\textsuperscript{2} & 218 cm\textsuperscript{2} & Does not fit \\
\bottomrule
\end{tabularx}
\end{table}

These are ideal signal floors, not net-flux predictions. With the 25-cm\textsuperscript{2}
aperture, the maximum ideal flows are approximately 2.29~L/min for $t_w=1$
and 0.69~L/min for $t_w=0.3$ at 10~A/m\textsuperscript{2}. Low-flow signal
acquisition and high-flow film characterisation should be separate modes.

\begin{figure}[htbp]
\centering
\includegraphics[width=0.94\textwidth]{../../outputs/eo_vapor_cell_sizing_v1/cell_dimension_tradeoffs.png}
\caption{Existing project calculation of cell dimensions and measurement trade-offs. Values are experiment-design sensitivities, not validated material performance.}
\end{figure}
\FloatBarrier

\section{Mechanical and fluidic starting point}

\begin{table}[htbp]
\centering
\caption{Fixture dimensions and their calculation basis.}
\begin{tabularx}{\textwidth}{@{}Y R Y@{}}
\toprule
Item & Starting value & Basis \\
\midrule
Maximum active aperture & 50 $\times$ 50 mm & Signal margin; low perimeter/area ratio \\
Physical membrane coupon & $\approx$80 $\times$ 80 mm & Double seals, guard groove and clamp land \\
Outer fixture & $\approx$110 $\times$ 110 mm & Ports, bolts and electrical isolation \\
Nominal gap per side & 2 mm & Local match to declared ERV channel \\
Gap validation & 1, 2 and 4 mm & Gas-film and pressure-drop sensitivity \\
Active volume per side & 5.00 mL & 25 cm\textsuperscript{2} $\times$ 0.2 cm \\
Headers, ports and taps & 1.222 mL & Proposed geometric allowance \\
Total swept volume & 6.222 mL/side & Active channel plus internal dead volume \\
Ideal residence at 6 L/min & 62 ms & 6.222 mL / 100 mL/s \\
Main pressure difference & within $\pm$5 Pa & Limit hydraulic transport artefacts \\
\bottomrule
\end{tabularx}
\end{table}

Hard stops, not gasket compression, must define the gas gap. For the 1-mm gap,
the plate-bow target is below 20~$\mu$m or the measured nonuniformity must be
included in analysis. The stiffness screen gives 15.35~mm PEEK, 5.81~mm
Al~6061-T6 and 4.16~mm 316L at a deliberately conservative 1-MPa equivalent
uniform load. These are screening values, not gasket-contact calculations;
final acceptance requires measured flatness under the real assembly load.

\section{Closed-cell fluctuation without through-flow}

At constant flux, the water shifted during one half-cycle is
\[
\Delta m_{\mathrm{half}}=JA\Delta t=\frac{JA}{2f},
\]
where $f$ is the full-cycle frequency. For 25~cm\textsuperscript{2} at the
208-g/m\textsuperscript{2}/h target, $JA=0.520$~g/h = 0.144~mg/s.

\begin{table}[htbp]
\centering
\caption{Ideal one-direction displacement for the 25-cm\textsuperscript{2} coupon.}
\begin{tabular}{@{}rrr@{}}
\toprule
Frequency & Half-cycle & Water shifted \\
\midrule
0.1 Hz & 5 s & 0.72 mg \\
0.01 Hz & 50 s & 7.2 mg \\
0.001 Hz & 500 s & 72 mg \\
0.0001 Hz & 5,000 s & 722 mg \\
\bottomrule
\end{tabular}
\end{table}

Multiply these values by 1.615 for the 336-g/m\textsuperscript{2}/h gross
reference. These linear values apply only while current and boundary states
remain available. They are initial command-rate calculations, not guaranteed
cycle amplitudes.

The 6.222-mL internal gas volume contains only about 0.1--0.2~mg water vapour
over the full 0--100\% RH span at 25--32~$^\circ$C. At 0.144~mg/s, the chamber
humidity can therefore change materially in roughly one second without an
external reservoir, sorbent or condensate capacity. The resulting activity
difference drives back-diffusion. Tubing, sensors, electrodes, supports and
walls must be included in the transient water inventory.

\begin{stopbox}
\textbf{Full-cycle result.} In the current constant-drag model, symmetric
charge-balanced polarity cycling gives
\[
\boxed{\text{net EO water transport over one complete cycle}=0.}
\]
Absolute shuttled water is nonzero, but the second half-cycle returns the EO
displacement. Continuous dehumidification requires a validated mechanism that
breaks this symmetry, such as flow switching, condensation, rectification,
asymmetric regeneration, pseudocapacitance or a defined Faradaic pathway.
\end{stopbox}

\section{Blocking-electrode capacity}

For purely blocking electrodes, the ideal finite displacement is
\[
m_{\max}=t_w\frac{M_w}{F}C_{\mathrm{pair},A}A\Delta V.
\]
The reference sensitivity uses 10,000~F/m\textsuperscript{2} for each identical
electrode, or 5,000~F/m\textsuperscript{2} for the series pair. For 25~cm\textsuperscript{2},
$t_w\approx100$ and a 1.8-V terminal limit, the transient model gives about
420~mg EO displacement before current collapses. The same ideal scaling at
$t_w=50$ gives approximately 210~mg. These capacitances and transport numbers
are unmeasured sensitivity coordinates.

\begin{figure}[htbp]
\centering
\includegraphics[width=0.88\textwidth]{../../outputs/capacitive_electrode_study_v1/figures/03_symmetric_ac_cycle.pdf}
\caption{Reference symmetric cycle from the existing transient study. Charge-balanced reversal closes the EO water balance over a full cycle.}
\end{figure}
\FloatBarrier

\section{Sensitivity map and experimental meaning}

\begin{table}[htbp]
\centering
\caption{First-order sensitivity and validity limits.}
\begin{tabularx}{\textwidth}{@{}Y Y Y@{}}
\toprule
Variable increased & First-order effect & Principal limitation \\
\midrule
Active area & Rate and capacitive inventory rise linearly & Edge leakage and current uniformity \\
Transport number & Gross flux per charge rises linearly & Must be measured in vapour \\
Current density & Initial gross flux rises linearly & Voltage, heating, polarisation, reactions \\
Half-cycle time & Ideal displacement rises linearly & Chamber/electrode saturation \\
Frequency & Displacement per half-cycle falls as $1/f$ & Fast cycling may probe electrodes, not membrane \\
Pair capacitance & Blocking-electrode inventory rises linearly & Still zero net in a symmetric full cycle \\
Gas flow & RH signal decreases at fixed water rate & Films and boundary control may improve \\
Channel gap & Changes volume, velocity, films and pressure drop & Must be measured in the 1/2/4-mm sweep \\
Diffusivity & Faster redistribution and passive flux & Back-diffusion can reduce net EO signal \\
\bottomrule
\end{tabularx}
\end{table}

The membrane redistribution scale is approximately $L^2/D$. For the
36.9-$\mu$m reference membrane, the present diffusivity screen spans seconds to
roughly one day. Steady state must therefore be defined by stable voltage and
closed donor/receiver water balances, not by one fixed elapsed time.

\section{Recommended test programme}

\begin{enumerate}
\item Run impermeable-foil, zero-current and passive-gradient blanks.
\item Use once-through gas at equal inlet activity and low current to locate a
measurable signal safely.
\item Apply both polarities in an ABBA sequence and calculate
$J_{\mathrm{odd}}=[J(+I)-J(-I)]/2$.
\item Require donor loss and receiver gain to close independently.
\item Repeat the 4/9/16/25-cm\textsuperscript{2} aperture sweep and fit area,
perimeter and blank terms.
\item Sweep flow separately from the 1/2/4-mm gap sweep to expose gas-film
limitations.
\item After sustained polarity-dependent transport is established, connect two
measured, well-mixed reservoirs and test reversible half-cycles.
\item Report initial rate, half-cycle displacement, peak-to-peak inventory,
full-cycle water closure, charge closure, voltage and energy.
\end{enumerate}

\begin{gobox}
\textbf{Proposed first success criterion.} Demonstrate repeatable
polarity-odd net transport of at least 200--250~g/m\textsuperscript{2}/h below
1~V under representative gas--gas conditions; show matching donor loss and
receiver gain; and show that the rate is dominated by active area rather than
perimeter or blank terms.
\end{gobox}

The resulting report should also give effective
$t_w=FJ_{\mathrm{odd}}/(jM_w)$, conductivity or area-specific resistance,
passive water flux, temperature, activity, pressure, flow, gap, electrode
state and measurement uncertainty for the same specimen.

\section{Conclusions and claim boundary}

The system model supplies a clear experimental scale: approximately
208~g/m\textsuperscript{2}/h additional EO-equivalent transport at O30, rounded
to a practical 200--250-g/m\textsuperscript{2}/h target. The 25-cm\textsuperscript{2}
coupon converts this into a measurable target rate of 0.520~g/h while retaining
smaller apertures for area/perimeter validation.

The fixture dimensions are driven by signal, boundary control, seal land,
channel similarity and artefact separation. They do not validate the assumed
high transport number. The decisive evidence must be a vapour-state,
polarity-dependent donor/receiver water balance under measured voltage and
current.

Closed-cell oscillation is useful but answers a different question. It measures
finite displacement, reversal and storage. Without through-flow or another
asymmetry, it cannot provide sustained net dehumidification. Absolute water
shuttled per half-cycle must never be reported as continuous moisture-removal
capacity.

\begin{warningbox}
All numerical material coordinates in this report remain conditional until
measured on the same membrane under representative vapour activity,
temperature, current density and interface conditions. The report defines an
experiment and falsifiable targets; it does not claim a validated membrane,
electrode, receiver or commercial ERV system.
\end{warningbox}

\section*{Reproduction and source documents}
\addcontentsline{toc}{section}{Reproduction and source documents}

The calculations are reproduced with:
\begin{verbatim}
python3 scripts/design_eo_vapor_test_cell.py
PYTHONPATH=src python3 scripts/run_capacitive_electrode_study_v1.py --overwrite
pytest -q -s tests/test_eo_vapor_test_cell_design.py
pytest -q -s tests/test_capacitive_electrode_study_v1.py
\end{verbatim}

The detailed bases are:
\begin{itemize}
\item \path{docs/EO_VAPOR_TEST_CELL_DESIGN.md};
\item \path{docs/TEST_CELL_AND_ERV_SCALE_CHECK_2026-09-01.md};
\item \path{docs/CAPACITIVE_ELECTRODE_TRANSIENT_STUDY_2026-09-01.md}; and
\item \path{docs/O30_REQUIRED_AREA_1V_v1.md}.
\end{itemize}

\end{document}
